\documentclass[10pt,a4paper]{amsart}
\usepackage{amsmath,amssymb,amsthm,graphicx}
% Retained minimal preamble for the Stacks Project chapter coherent.tex (Stacks chapter, pre-delegation retained fragment for stacks-coherent).
% Verbatim excerpt of the chapter preamble (sources/stacks-spaces-descent/preamble.tex): the reference/slogan/history
% environment block (native lines 7-12), Xy-pic setup (native lines 14-19) and the native theorem declarations with
% the mathematical macros (native lines 154-217). Copied mechanically; nothing added, removed or redefined. The
% wrapper class and the amsmath/amssymb/amsthm/graphicx base packages are supplied by the assembler, so the
% class/xr-hyper/hyperref/externaldocument part of the native preamble is omitted.
% For dealing with references we use the comment environment
\usepackage{verbatim}
\newenvironment{reference}{\comment}{\endcomment}
%\newenvironment{reference}{}{}
\newenvironment{slogan}{\comment}{\endcomment}
\newenvironment{history}{\comment}{\endcomment}
% For commutative diagrams we use Xy-pic
\usepackage[all]{xy}

% We use 2cell for 2-commutative diagrams.
\xyoption{2cell}
\UseAllTwocells
% Theorem environments.
%
\theoremstyle{plain}
\newtheorem{theorem}[subsection]{Theorem}
\newtheorem{proposition}[subsection]{Proposition}
\newtheorem{lemma}[subsection]{Lemma}

\theoremstyle{definition}
\newtheorem{definition}[subsection]{Definition}
\newtheorem{example}[subsection]{Example}
\newtheorem{exercise}[subsection]{Exercise}
\newtheorem{situation}[subsection]{Situation}

\theoremstyle{remark}
\newtheorem{remark}[subsection]{Remark}
\newtheorem{remarks}[subsection]{Remarks}

\numberwithin{equation}{subsection}

% Macros
%
\def\lim{\mathop{\mathrm{lim}}\nolimits}
\def\colim{\mathop{\mathrm{colim}}\nolimits}
\def\Spec{\mathop{\mathrm{Spec}}}
\def\Hom{\mathop{\mathrm{Hom}}\nolimits}
\def\Ext{\mathop{\mathrm{Ext}}\nolimits}
\def\SheafHom{\mathop{\mathcal{H}\!\mathit{om}}\nolimits}
\def\SheafExt{\mathop{\mathcal{E}\!\mathit{xt}}\nolimits}
\def\Sch{\mathit{Sch}}
\def\Mor{\mathop{\mathrm{Mor}}\nolimits}
\def\Ob{\mathop{\mathrm{Ob}}\nolimits}
\def\Sh{\mathop{\mathit{Sh}}\nolimits}
\def\NL{\mathop{N\!L}\nolimits}
\def\CH{\mathop{\mathrm{CH}}\nolimits}
\def\proetale{{pro\text{-}\acute{e}tale}}
\def\etale{{\acute{e}tale}}
\def\QCoh{\mathit{QCoh}}
\def\Ker{\mathop{\mathrm{Ker}}}
\def\Im{\mathop{\mathrm{Im}}}
\def\Coker{\mathop{\mathrm{Coker}}}
\def\Coim{\mathop{\mathrm{Coim}}}

% Boxtimes
%
\DeclareMathSymbol{\boxtimes}{\mathbin}{AMSa}{"02}

%
% Macros for moduli stacks/spaces
%
\def\QCohstack{\mathcal{QC}\!\mathit{oh}}
\def\Cohstack{\mathcal{C}\!\mathit{oh}}
\def\Spacesstack{\mathcal{S}\!\mathit{paces}}
\def\Quotfunctor{\mathrm{Quot}}
\def\Hilbfunctor{\mathrm{Hilb}}
\def\Curvesstack{\mathcal{C}\!\mathit{urves}}
\def\Polarizedstack{\mathcal{P}\!\mathit{olarized}}
\def\Complexesstack{\mathcal{C}\!\mathit{omplexes}}
% \Pic is the operator that assigns to X its picard group, usage \Pic(X)
% \Picardstack_{X/B} denotes the Picard stack of X over B
% \Picardfunctor_{X/B} denotes the Picard functor of X over B
\def\Pic{\mathop{\mathrm{Pic}}\nolimits}
\def\Picardstack{\mathcal{P}\!\mathit{ic}}
\def\Picardfunctor{\mathrm{Pic}}
\def\Deformationcategory{\mathcal{D}\!\mathit{ef}}
\usepackage{cleveref}
\crefname{originallemma}{original lemma}{original lemmas}
\Crefname{originallemma}{Original lemma}{Original lemmas}
\makeatletter
\newlabel{algebra-lemma-completion-flat}{{00MB}{source}{}{sourcepointer-1}{}}
\newlabel{algebra-lemma-completion-flat@cref}{{[originallemma][0][]00MB}{[1][0][]source}{}{}{}}
\makeatother
\crefname{originallemma}{original lemma}{original lemmas}
\Crefname{originallemma}{Original lemma}{Original lemmas}
\makeatletter
\newlabel{lemma-ML-cohomology-powers-ideal}{{02OB}{source}{}{sourcepointer-2}{}}
\newlabel{lemma-ML-cohomology-powers-ideal@cref}{{[originallemma][0][]02OB}{[1][0][]source}{}{}{}}
\makeatother
\crefname{originallemma}{original lemma}{original lemmas}
\Crefname{originallemma}{Original lemma}{Original lemmas}
\makeatletter
\newlabel{lemma-flat-base-change-cohomology}{{02KH}{source}{}{sourcepointer-3}{}}
\newlabel{lemma-flat-base-change-cohomology@cref}{{[originallemma][0][]02KH}{[1][0][]source}{}{}{}}
\makeatother
\crefname{originallemma}{original lemma}{original lemmas}
\Crefname{originallemma}{Original lemma}{Original lemmas}
\makeatletter
\newlabel{lemma-fully-faithful}{{0883}{source}{}{sourcepointer-4}{}}
\newlabel{lemma-fully-faithful@cref}{{[originallemma][0][]0883}{[1][0][]source}{}{}{}}
\makeatother
\crefname{originallemma}{original lemma}{original lemmas}
\Crefname{originallemma}{Original lemma}{Original lemmas}
\makeatletter
\newlabel{lemma-inverse-systems-affine}{{087W}{source}{}{sourcepointer-5}{}}
\newlabel{lemma-inverse-systems-affine@cref}{{[originallemma][0][]087W}{[1][0][]source}{}{}{}}
\makeatother
\crefname{originallemma}{original lemma}{original lemmas}
\Crefname{originallemma}{Original lemma}{Original lemmas}
\makeatletter
\newlabel{lemma-power-ideal-kills-sheaf}{{01Y9}{source}{}{sourcepointer-6}{}}
\newlabel{lemma-power-ideal-kills-sheaf@cref}{{[originallemma][0][]01Y9}{[1][0][]source}{}{}{}}
\makeatother
\crefname{originallemma}{original lemma}{original lemmas}
\Crefname{originallemma}{Original lemma}{Original lemmas}
\makeatletter
\newlabel{more-algebra-lemma-neighbourhood-equivalence}{{05EC}{source}{}{sourcepointer-7}{}}
\newlabel{more-algebra-lemma-neighbourhood-equivalence@cref}{{[originallemma][0][]05EC}{[1][0][]source}{}{}{}}
\makeatother
\crefname{originallemma}{original lemma}{original lemmas}
\Crefname{originallemma}{Original lemma}{Original lemmas}
\makeatletter
\newlabel{proposition-proper-pushforward-coherent}{{02O5}{source}{}{sourcepointer-8}{}}
\newlabel{proposition-proper-pushforward-coherent@cref}{{[originallemma][0][]02O5}{[1][0][]source}{}{}{}}
\makeatother
\crefname{originallemma}{original lemma}{original lemmas}
\Crefname{originallemma}{Original lemma}{Original lemmas}
\makeatletter
\newlabel{remark-inverse-systems-kernel-cokernel-annihilated-by}{{0888}{source}{}{sourcepointer-9}{}}
\newlabel{remark-inverse-systems-kernel-cokernel-annihilated-by@cref}{{[originallemma][0][]0888}{[1][0][]source}{}{}{}}
\makeatother
\crefname{originallemma}{original lemma}{original lemmas}
\Crefname{originallemma}{Original lemma}{Original lemmas}
\makeatletter
\newlabel{theorem-formal-functions}{{02OC}{source}{}{sourcepointer-10}{}}
\newlabel{theorem-formal-functions@cref}{{[originallemma][0][]02OC}{[1][0][]source}{}{}{}}
\makeatother
\title{Push-pull for coherent systems under a proper map that is an isomorphism off V(K): for Y Noetherian, J, K quasi-coherent ideal sheaves, f : X -> Y proper and an isomorphism over V = Y \textbackslash{} V(K), I = f\textasciicircum{}\{-1\}J O\_X, (G\_n) in Coh(Y, J), F coherent on X and beta : (f\textasciicircum{}*G\_n) -> F\textasciicircum{} an isomorphism in Coh(X, I), there exists alpha : (G\_n) -> (f\_*F)\textasciicircum{} in Coh(Y, J) whose kernel and cokernel are annihilated by a power of K.}
\author{The Stacks Project Authors}
\date{stacks-project:coherent.tex:local-snapshot-20261012; GNU FDL 1.2 or later}
\begin{document}
\maketitle
\noindent\textit{Editorial note.} Statement span is lines 7241-7261 of sources/stacks-coherent/pinned-original-source.tex; proof span is lines 7263-7357, the complete author proof (93 body lines).\par\medskip
\noindent\textit{Editorial note.} Navigation stubs: 10 cited-but-not-included labels are supplied as non-load-bearing pointers with their pinned tags where the snapshot tag index has them.\par\medskip
\noindent\textit{Editorial note.} extra\_packages is empty on purpose: the retained preamble already loads xy with options and a second declaration clashes in the assembler wrapper.\par\medskip
\noindent\textit{Editorial note.} Text is retained verbatim, including the author's own typos and placeholders.\par\medskip
\noindent\textit{Editorial note.} Grothendieck existence theorem machinery; full derivation in span, no omitted steps beyond pointers to cited lemmas used as inputs.\par\medskip
\noindent\textit{Editorial note.} Span computed from the \textbackslash{}label marker through the enclosing \textbackslash{}end\{lemma\}; precard\_gate RESULT=OK (probe L9999).\par\medskip
\noindent\textit{Editorial note.} Cross-chapter same-label warning (NOT an occupancy conflict, Chief ruling 2026-10-13 evening): label 'lemma-inverse-systems-push-pull' also exists in stacks-spaces-more-morphisms — a different proposition in a different chapter file.\par\medskip
\noindent\textit{Allowed standard prerequisites (editorial).}\begin{itemize}
\item Explicit prerequisite boundary. See why-field; every pointer named here occurs verbatim in the retained statement/proof text.
\end{itemize}\medskip
\ifdefined\hypertarget\hypertarget{sourcepointer-1}{}\fi
\noindent\textit{Source reference.} Original lemma 00MB: sources/stacks-coherent/pinned-original-source.tex, source lines 7357--7357. Cited by the retained statement or proof but not part of the collected spans; supplied only as a navigation pointer. The argument does not depend on its content.\par\medskip
\ifdefined\hypertarget\hypertarget{sourcepointer-2}{}\fi
\noindent\textit{Source reference.} Original lemma 02OB: sources/stacks-coherent/pinned-original-source.tex, source lines 5770--5797. Cited by the retained statement or proof but not part of the collected spans; supplied only as a navigation pointer. The argument does not depend on its content.\par\medskip
\ifdefined\hypertarget\hypertarget{sourcepointer-3}{}\fi
\noindent\textit{Source reference.} Original lemma 02KH: sources/stacks-coherent/pinned-original-source.tex, source lines 1165--1186. Cited by the retained statement or proof but not part of the collected spans; supplied only as a navigation pointer. The argument does not depend on its content.\par\medskip
\ifdefined\hypertarget\hypertarget{sourcepointer-4}{}\fi
\noindent\textit{Source reference.} Original lemma 0883: sources/stacks-coherent/pinned-original-source.tex, source lines 6977--6982. Cited by the retained statement or proof but not part of the collected spans; supplied only as a navigation pointer. The argument does not depend on its content.\par\medskip
\ifdefined\hypertarget\hypertarget{sourcepointer-5}{}\fi
\noindent\textit{Source reference.} Original lemma 087W: sources/stacks-coherent/pinned-original-source.tex, source lines 6489--6495. Cited by the retained statement or proof but not part of the collected spans; supplied only as a navigation pointer. The argument does not depend on its content.\par\medskip
\ifdefined\hypertarget\hypertarget{sourcepointer-6}{}\fi
\noindent\textit{Source reference.} Original lemma 01Y9: sources/stacks-coherent/pinned-original-source.tex, source lines 2786--2793. Cited by the retained statement or proof but not part of the collected spans; supplied only as a navigation pointer. The argument does not depend on its content.\par\medskip
\ifdefined\hypertarget\hypertarget{sourcepointer-7}{}\fi
\noindent\textit{Source reference.} Original lemma 05EC: sources/stacks-coherent/pinned-original-source.tex, source lines 7357--7357. Cited by the retained statement or proof but not part of the collected spans; supplied only as a navigation pointer. The argument does not depend on its content.\par\medskip
\ifdefined\hypertarget\hypertarget{sourcepointer-8}{}\fi
\noindent\textit{Source reference.} Original lemma 02O5: sources/stacks-coherent/pinned-original-source.tex, source lines 5325--5328. Cited by the retained statement or proof but not part of the collected spans; supplied only as a navigation pointer. The argument does not depend on its content.\par\medskip
\ifdefined\hypertarget\hypertarget{sourcepointer-9}{}\fi
\noindent\textit{Source reference.} Original lemma 0888: sources/stacks-coherent/pinned-original-source.tex, source lines 7104--7124. Cited by the retained statement or proof but not part of the collected spans; supplied only as a navigation pointer. The argument does not depend on its content.\par\medskip
\ifdefined\hypertarget\hypertarget{sourcepointer-10}{}\fi
\noindent\textit{Source reference.} Original lemma 02OC: sources/stacks-coherent/pinned-original-source.tex, source lines 5865--5887. Cited by the retained statement or proof but not part of the collected spans; supplied only as a navigation pointer. The argument does not depend on its content.\par\medskip
\begin{lemma}
\label{lemma-inverse-systems-push-pull}
Let $Y$ be a Noetherian scheme. Let
$\mathcal{J}, \mathcal{K} \subset \mathcal{O}_Y$
be quasi-coherent sheaves of ideals.
Let $f : X \to Y$ be a proper morphism which is an isomorphism
over $V = Y \setminus V(\mathcal{K})$.
Set $\mathcal{I} = f^{-1}\mathcal{J} \mathcal{O}_X$.
Let $(\mathcal{G}_n)$ be an object of $\textit{Coh}(Y, \mathcal{J})$,
let $\mathcal{F}$ be a coherent $\mathcal{O}_X$-module, and let
$\beta : (f^*\mathcal{G}_n)  \to \mathcal{F}^\wedge$ be an isomorphism in
$\textit{Coh}(X, \mathcal{I})$. Then there exists a map
$$
\alpha :
(\mathcal{G}_n)
\longrightarrow
(f_*\mathcal{F})^\wedge
$$
in $\textit{Coh}(Y, \mathcal{J})$ whose kernel and cokernel
are annihilated by a power of $\mathcal{K}$.
\end{lemma}

\begin{proof}
Since $f$ is a proper morphism we see that $f_*\mathcal{F}$
is a coherent $\mathcal{O}_Y$-module
(Proposition \ref{proposition-proper-pushforward-coherent}).
Thus the statement of the lemma makes sense.
Consider the compositions
$$
\gamma_n : \mathcal{G}_n \to
f_*f^*\mathcal{G}_n \to
f_*(\mathcal{F}/\mathcal{I}^n\mathcal{F}).
$$
Here the first map is the adjunction map and the second is $f_*\beta_n$.
We claim that there exists a unique $\alpha$ as in the lemma
such that the compositions
$$
\mathcal{G}_n \xrightarrow{\alpha_n}
f_*\mathcal{F}/\mathcal{J}^nf_*\mathcal{F} \to
f_*(\mathcal{F}/\mathcal{I}^n\mathcal{F})
$$
equal $\gamma_n$ for all $n$. Because of the uniqueness we may assume
that $Y = \Spec(B)$ is affine. Let $J \subset B$ corresponds to the
ideal $\mathcal{J}$. Set
$$
M_n = H^0(X, \mathcal{F}/\mathcal{I}^n\mathcal{F})
\quad\text{and}\quad
M = H^0(X, \mathcal{F})
$$
By Lemma \ref{lemma-ML-cohomology-powers-ideal} and
Theorem \ref{theorem-formal-functions}
the inverse limit of the modules
$M_n$ equals the completion $M^\wedge = \lim M/J^nM$.
Set $N_n = H^0(Y, \mathcal{G}_n)$ and $N = \lim N_n$.
Via the equivalence of categories of
Lemma \ref{lemma-inverse-systems-affine}
the finite $B^\wedge$ modules $N$ and $M^\wedge$ correspond
to $(\mathcal{G}_n)$ and $f_*\mathcal{F}^\wedge$.
It follows from this that $\alpha$ has to be the morphism of
$\textit{Coh}(Y, \mathcal{J})$ corresponding to the homomorphism
$$
\lim \gamma_n : N = \lim_n N_n \longrightarrow \lim M_n = M^\wedge
$$
of finite $B^\wedge$-modules.

\medskip\noindent
We still have to show that the kernel and cokernel of $\alpha$ are
annihilated by a power of $\mathcal{K}$. Set $Y' = \Spec(B^\wedge)$
and $X' = Y' \times_Y X$. Let $\mathcal{K}'$, $\mathcal{J}'$, $\mathcal{G}'_n$
and $\mathcal{I}'$, $\mathcal{F}'$ be the pullback of
$\mathcal{K}$, $\mathcal{J}$, $\mathcal{G}_n$ and
$\mathcal{I}$, $\mathcal{F}$, to $Y'$ and $X'$.
The projection morphism $f' : X' \to Y'$ is the base change of
$f$ by $Y' \to Y$. Note that $Y' \to Y$ is a flat morphism of schemes
as $B \to B^\wedge$ is flat by
Algebra, Lemma \ref{algebra-lemma-completion-flat}.
Hence $f'_*\mathcal{F}'$, resp.\ $f'_*(f')^*\mathcal{G}_n'$
is the pullback of $f_*\mathcal{F}$, resp.\ $f_*f^*\mathcal{G}_n$
to $Y'$ by Lemma \ref{lemma-flat-base-change-cohomology}.
The uniqueness of our construction shows the pullback of $\alpha$ to $Y'$
is the corresponding map $\alpha'$ constructed for the situation on $Y'$.
Moreover, to check that the kernel and cokernel of $\alpha$ are
annihilated by $\mathcal{K}^t$ it suffices to check that the
kernel and cokernel of $\alpha'$ are annihilated by
$(\mathcal{K}')^t$. Namely, to see this we need to check this for
kernels and cokernels of the maps $\alpha_n$ and $\alpha'_n$
(see Remark \ref{remark-inverse-systems-kernel-cokernel-annihilated-by})
and the ring map $B \to B^\wedge$ induces
an equivalence of categories between modules annihilated by
$J^n$ and $(J')^n$, see
More on Algebra, Lemma \ref{more-algebra-lemma-neighbourhood-equivalence}.
Thus we may assume $B$ is complete with respect to $J$.

\medskip\noindent
Assume $Y = \Spec(B)$ is affine, $\mathcal{J}$ corresponds to the ideal
$J \subset B$, and $B$ is complete with respect to $J$.
In this case $(\mathcal{G}_n)$ is in the essential image of the functor
$\textit{Coh}(\mathcal{O}_Y) \to \textit{Coh}(Y, \mathcal{J})$.
Say $\mathcal{G}$ is a coherent $\mathcal{O}_Y$-module such that
$(\mathcal{G}_n) = \mathcal{G}^\wedge$. Note that
$f^*(\mathcal{G}^\wedge) = (f^*\mathcal{G})^\wedge$. Hence
Lemma \ref{lemma-fully-faithful}
tells us that $\beta$ comes from an isomorphism
$b : f^*\mathcal{G} \to \mathcal{F}$
and $\alpha$ is the completion functor applied to
$$
\mathcal{G} \to f_*f^*\mathcal{G} \cong f_*\mathcal{F}
$$
Hence we are trying to verify that the kernel and cokernel of the
adjunction map $c : \mathcal{G} \to f_*f^*\mathcal{G}$ are annihilated by
a power of $\mathcal{K}$. However, since the restriction
$f|_{f^{-1}(V)} : f^{-1}(V) \to V$ is an isomorphism
we see that $c|_V$ is an isomorphism. Thus the coherent sheaves
$\Ker(c)$ and $\Coker(c)$ are supported on $V(\mathcal{K})$
hence are annihilated by a power of $\mathcal{K}$
(Lemma \ref{lemma-power-ideal-kills-sheaf}) as desired.
\end{proof}

\end{document}
